\subsection{Answer Set Programming}
\label{sec:foundations_asp}

Dimopoulos et al.~\cite{Dimopoulos1997} describe \ac{ASP} as a declarative programming paradigm based on stable model semantics for logic programs. This section provides formal definitions of \ac{ASP} concepts used in the implementation section.

\begin{defn}[Logic Program~\cite{Dimopoulos1997}]
    \label{defn:logic_program}
    A \emph{logic program} $\Pi$ is a finite set of rules. Each rule has the form:
    \begin{equation*}
        h \leftarrow b_1, \ldots, b_m, \mathbf{not}\ c_1, \ldots, \mathbf{not}\ c_n
    \end{equation*}
    where $h$ is the \emph{head} (an atom or empty), $b_1, \ldots, b_m$ are the \emph{positive body} atoms, and $c_1, \ldots, c_n$ are the \emph{negative body} atoms preceded by default negation $\mathbf{not}$.
\end{defn}

A rule with an empty head is a \emph{constraint} (denoted $\leftarrow b_1, \ldots, b_m$), which eliminates answer sets containing the body atoms. A rule with an empty body is a \emph{fact} (denoted $h.$).

\begin{defn}[Answer Set~\cite{Dimopoulos1997}]
    \label{defn:answer_set}
    An \emph{answer set} (or \emph{stable model}) of a logic program $\Pi$ is a minimal set of atoms $S$ such that:
    \begin{enumerate}
        \item For every rule $h \leftarrow b_1, \ldots, b_m, \mathbf{not}\ c_1, \ldots, \mathbf{not}\ c_n$ in $\Pi$: if $\{b_1, \ldots, b_m\} \subseteq S$ and $\{c_1, \ldots, c_n\} \cap S = \emptyset$, then $h \in S$.
        \item $S$ is minimal with respect to set inclusion among all sets satisfying condition 1.
    \end{enumerate}
\end{defn}

\begin{example}[Simple ASP Program]
    Consider the program with rules: $a.$ (fact), $b \leftarrow a.$ (conditional), $\leftarrow a, c.$ (constraint). The unique answer set is $\{a, b\}$: the fact establishes $a$, the conditional rule derives $b$ from $a$, and $c$ is not derived so the constraint is not violated.
\end{example}

Eiter and Gottlob~\cite{Eiter1995} formalize two complementary modes of reasoning over the answer sets of a logic program.

\begin{defn}[Brave and Cautious Consequences~\cite{Eiter1995}]
    \label{defn:brave_cautious}
    Let $\Pi$ be a logic program with answer sets $\text{AS}(\Pi) = \{S_1, \ldots, S_n\}$. The \emph{brave consequences} and \emph{cautious consequences} of $\Pi$ are defined as:
    \begin{align*}
        \text{BRAVE}(\Pi)    & = \bigcup_{S \in \text{AS}(\Pi)} S &  & \text{(atoms true in at least one answer set)} \\
        \text{CAUTIOUS}(\Pi) & = \bigcap_{S \in \text{AS}(\Pi)} S &  & \text{(atoms true in every answer set)}
    \end{align*}
\end{defn}

Modern \ac{ASP} solvers such as Clingo, developed by Gebser et al.~\cite{Gebser2012}, support both modes natively. In a planning context with non-deterministic operator selection, each answer set corresponds to a different execution trace. Brave consequences capture everything achievable under any choice of operators, making brave reasoning the appropriate mode for reachability analysis.

\ac{ASP} provides a natural framework for encoding planning problems. The modern tool plasp~3 by Dimopoulos et al.~\cite{Dimopoulos2019} translates \ac{PDDL} specifications into \ac{ASP} programs where answer sets correspond to valid plans. The system generates rules encoding operator applicability, effect application, and goal satisfaction.

\begin{defn}[Sequential Planning Encoding~\cite{Dimopoulos2019}]
    \label{defn:sequential_encoding}
    A \emph{sequential encoding} for horizon $i$, denoted $S(i)$, is an \ac{ASP} program that searches for plans of exactly $i$ steps. The encoding includes time-indexed fluents $\text{holds}(p, t)$ for propositions $p$ at time steps $t \in \{0, 1, \ldots, i\}$, action atoms $\text{occurs}(o, t)$ indicating operator $o$ executes at step $t$, and constraints ensuring exactly one action occurs at each step.
\end{defn}

Sequential encodings enforce that exactly one operator executes at each time step. To find plans of unknown length, the encoding is solved iteratively for increasing horizons $i = 0, 1, 2, \ldots$ until either an answer set is found or a termination condition is met.

Parallel encodings relax the single-action restriction, allowing multiple non-interfering actions to execute in parallel at each step. To ensure valid parallel execution, these encodings include \emph{serialization constraints} that prevent conflicting actions from occurring at the same time step. Two actions conflict if one deletes a precondition or add effect of the other, or if both add and delete the same proposition. Serialization constraints take the form $\leftarrow \text{occurs}(o_1, t), \text{occurs}(o_2, t)$ for each pair of conflicting actions $o_1$ and $o_2$.

For inconsistency measurement applications, \ac{ASP} offers several advantages. First, the declarative nature of \ac{ASP} programs allows for systematic analysis of conflicts through unsatisfiable cores and related computational structures, as demonstrated by Ulbricht et al.~\cite{Ulbricht2016} in their adaptation of inconsistency measures to \ac{ASP}. Second, \ac{ASP} provides natural support for non-monotonic reasoning through default negation, which is essential for planning domains where the monotonicity postulate of classical inconsistency measures may not hold due to the frame problem.

Regarding computational performance, the choice between \ac{ASP} and \ac{SAT}-based approaches depends on the specific measures and encodings. Kuhlmann et al.~\cite{Kuhlmann2025} found that \ac{ASP}-based algorithms outperform basic \ac{SAT} encodings for several inconsistency measures. However, Niskanen et al.~\cite{Niskanen2023} demonstrated that specialized MaxSAT approaches can significantly outperform \ac{ASP} methods for the same measures. This thesis employs \ac{ASP} for implementation due to its expressiveness for encoding reachability-based measures and its natural support for the non-monotonic reasoning patterns that arise in planning diagnostics.
